\pdfsuppresswarningpagegroup=1
\documentclass{article}

\usepackage[margin=1in]{geometry}
\usepackage{newtxtext,newtxmath}
\usepackage{amsmath}

\usepackage[
  pdfauthor={William C. Dawn},
  pdftitle={Analytic Solutions of the Heat Conduction Equation for Validation},
  pdfcreator={pdflatex},
  pdfsubject={Heat Conduction Validation},
  pdfkeywords={},
  colorlinks=true,
  linkcolor=black,
  citecolor=black,
  filecolor=black,
  urlcolor=black,
]{hyperref}

\usepackage{graphicx,subcaption}
\usepackage[section]{placeins}
\usepackage{booktabs}

\title{Analytic Solutions of the Heat Conduction Equation for Validation}
\author{William C. Dawn}

\newcommand{\grad}{\nabla}

\newcommand{\vr}{\mathbf{r}}
\newcommand{\kbar}{\bar{k}}
\newcommand{\that}{\hat{T}}

% reference macros
\newcommand{\eref}[1]{\hyperref[#1]{Eq.~}(\ref{#1})}
\newcommand{\fref}[1]{\hyperref[#1]{Fig.~}\ref{#1}}
\newcommand{\tref}[1]{\hyperref[#1]{Table~}\ref{#1}}
\newcommand{\sref}[1]{\hyperref[#1]{\S}\ref{#1}}
\newcommand{\algorithmref}[1]{\hyperref[#1]{Algorithm~}\ref{#1}}

\begin{document}

\maketitle

\section{Purpose}

This note provides analytic solutions to the heat conduction equation that can be used to validate a computer program.
These solutions were originally developed for use with microreactor modeling so special attention is paid to solutions with heat pipes.

\section{General Solution}

The general form of the steady-state heat conduction equation is
\begin{equation}
  \label{eq:general_heat_conduction}
  - \grad \cdot ( k(\vr, T(\vr)) \, \grad T(\vr) ) = q(\vr)
\end{equation}
where $k$ is the thermal conductivity of the material, $T$ is the temperature distribution, $q$ is the volumetric heat source, and $\vr$ is the position vector.

Note that in its most general form, the thermal conductivity depends on both position (i.e. the material at that position) and temperature.
Physically speaking, thermal conductivity is temperature dependent.
However, \eref{eq:general_heat_conduction} can be solved using an fixed-point iteration algorithm so solvers typically only need to consider thermal conductivity as a piecewise-constant function of position.
The fixed-point iteration algorithm is not guaranteed to converge as it would require knowledge of the properties of the thermal conductivity function, but this method is typically acceptable for physical simulations.

When evaluating integrals to solve \eref{eq:general_heat_conduction}, one must appropriately consider the Laplacian in the space (e.g. for cylindrical coordinates).
Specific geometries and heat source functions are considered next.

\section{Heat Pipe Modeling}

Several microreactor designs are considering the use of heat pipes for cooling.
Heat pipes are thin cylindrical tubes filled with a mesh/wick material and a small amount of a coolant (e.g. liquid sodium) such that the heat pipe is only partially filled with fluid.
The heat removal of a heat pipe is driven by the evaporation of the fluid and capillary action and therefore these can be used without gravity (in space) and do not require the modelling of a moving fluid.

Any attempt to model the temperatures or heat flow within a heat pipe is beyond the scope of this work.
For the purposes of this model, heat pipes are assumed to operate at a constant temperature.
If a more accurate heat pipe model is required, an appropriate model could be added to the computer program and solved as an additional multiphysics model.

\section{One-Dimension Cartesian Heat Conduction}
\label{sec:oned}

Begin by considering the heat conduction equation in one-dimension
\begin{equation}
  \label{eq:oned_heat_conduction}
  -k \; \frac{d^2}{dx^2} T(x) = q(x).
\end{equation}

For this specific problem, the geometry in \fref{fig:linear_geometry} will be used.
\fref{fig:linear_geometry} shows a one-dimension slab of length $L$ with a fixed temperature boundary condition $T(0)=T_{HP}$ (constant heat pipe temperature) and a fixed temperature boundary condition $T(L)=T_{B}$ (constant external boundary temperature).
The heat source function considered is
\begin{equation}
  \label{eq:oned_q}
  q(x) = q_0 \sin \left( \frac{\pi \, x}{L} \right)
\end{equation}
where $q_0$ is a scaling constant.

\begin{figure}
  \centering
  \includegraphics[width=0.7\textwidth]{../figs/linear_geometry}
  \caption{One-Dimension Cartesian Heat Conduction Geometry.}
  \label{fig:linear_geometry}
\end{figure}

Integrating \eref{eq:oned_heat_conduction} twice yields
\begin{equation}
  \label{eq:oned_general_solution}
  T(x) = \frac{q_0}{k} \left( \frac{L}{\pi} \right)^2 \left( \sin \left( \frac{\pi \, x}{L} \right) + c_1 x + c_2 \right)
\end{equation}
where $c_1$ and $c_2$ are constants of integration whose values will be obtained by imposing the boundary conditions.
Beginning with the boundary condition at $x=0$.
\begin{align}
  T(0) &= T_{HP} \\
  &= \frac{q_0}{k} \left( \frac{L}{\pi} \right)^2 c_2  \\
  \label{eq:oned_c2}
  \therefore c_2 &= T_{HP} \frac{k}{q_0} \left( \frac{\pi}{L} \right)^2
\end{align}
Substituting \eref{eq:oned_c2} into \eref{eq:oned_general_solution}.
\begin{equation}
  \label{eq:oned_intermediate_solution}
  T(x) = \frac{q_0}{k} \left( \frac{L}{\pi} \right)^2 \left( \sin \left( \frac{\pi \, x}{L} \right) + c_1 x \right) + T_{HP}
\end{equation}
Next, evaluate the boundary condition at $x=L$.
\begin{align}
  T(L) &= T_B \\
  &= \frac{q_0}{k} \left( \frac{L}{\pi} \right)^2 c_1 x + T_{HP} \\
  \label{eq:oned_c1}
  \therefore c_1 &= \frac{k}{q_0} \left( \frac{L}{\pi} \right)^2 \frac{T_B - T_{HP}}{L}
\end{align}
Finally, substitute \eref{eq:oned_c1} into \eref{eq:oned_intermediate_solution} to obtain the final form of the solution.
\begin{equation}
  \label{eq:oned_solution}
  T(x) = \frac{q_0}{k} \left( \frac{L}{\pi} \right)^2 \sin \left( \frac{\pi \, x}{L} \right) + \frac{T_B - T_{HP}}{L} x + T_{HP}
\end{equation}
Using the (dimensionless) values in \tref{tab:oned}, an example solution is plotted in \fref{fig:linear}.

\begin{table}
  \centering
  \caption{One-Dimension Cartesian Heat Conduction Constants.}
  \label{tab:oned}
  \begin{tabular}{cc}
    \\
    \toprule
    Variable & Value \\
    \midrule
    $T_{HP}$ & 400 \\
    $T_{B}$  & 300 \\
    $L$      & 1   \\
    $k$      & 0.1 \\
    $q_0$    & 100 \\
    \bottomrule
  \end{tabular}
\end{table}

\begin{figure}
  \centering
  \includegraphics[width=0.7\textwidth]{../figs/linear}
  \caption{One-Dimension Cartesian Heat Conduction Solution.}
  \label{fig:linear}
\end{figure}

\section{One-Dimension Cylindrical Heat Conduction with Heat Pipe}

Begin by considering the heat conduction equation in one-dimension cylindrical geometry.
\begin{equation}
  \label{eq:cyl_heat_conduction}
  -k \; \frac{d}{dr} \left( \frac{1}{r} \frac{d}{dr} \left( r \, T(r) \right) \right) = q(r)
\end{equation}
Note the proper treatment of the Laplacian for one-dimension cylindrical geometry in \eref{eq:cyl_heat_conduction}.

This problem is designed to model the treatment of heat pipes in a thermal conduction solver.
The geometry for this problem is shown in \fref{fig:radial_geometry}.
The model is a cylindrical heat pipe surrounded by a heat generating and conducting medium with thermal conductivity $k$ and heat source $q(r)$.
When the heat pipe is treated as a constant temperature boundary condition, this is a one-dimension cylindrical model with boundary conditions $T(R_{HP}) = T_{HP}$ and $T(R_B) = T_B$ where $R_{HP}$ is the radius of the heat pipe and $R_B$ is the radius of the external boundary.
The heat source function considered is
\begin{equation}
  \label{eq:cyl_source}
  q(r) = q_0 \frac{R_B - r}{R_B - R_{HP}}
\end{equation}
where $q_0$ is a scaling constant.
\eref{eq:cyl_source} has the property that it is linear and $q(R_{HP})=q_0$ and $q(R_B)=0$. 

\begin{figure}
  \centering
  \includegraphics[width=0.5\textwidth]{../figs/radial_geometry}
  \caption{One-Dimension Cylindrical Heat Conduction Geometry.}
  \label{fig:radial_geometry}
\end{figure}

Integrating \eref{eq:cyl_heat_conduction} twice in cylindrical geometry yields
\begin{equation}
  \label{eq:cyl_general}
  T(r) = \frac{q_0}{k} \frac{1}{R_B-R_{HP}} \left( \frac{r^3}{9} - \frac{R_B \, r^2}{4} + c_1 \ln(r) + c_2 \right)
\end{equation}
where $c_1$ and $c_2$ are integration constants that will be obtained by evaluating the boundary conditions.

A warning: there is very little cancellation in this problem and lots of tedious algebra.
Begin by evaluating the heat pipe boundary condition at $r=R_{HP}$.
\begin{align}
  T(R_{HP}) &= T_{HP} \\
  &= \frac{q_0}{k} \frac{1}{R_B - R_{HP}} \left( \frac{R_{HP}^3}{9} - \frac{R_{B} \, R_{HP}^2}{4} + c_1 \ln(R_{HP}) + c2 \right) \\
  \label{eq:cyl_c2_expression1}
  \therefore c_2 &= T_{HP} \frac{k}{q_0} (R_B - R_{HP}) - \frac{R_{HP}^3}{9} + \frac{R_B \, R_{HP}^2}{4} - c_1 \ln(R_{HP})
\end{align}
Continue by evaluating the external boundary condition at $r=R_B$.
\begin{align}
  T(R_B) &= T_B \\
  &= \frac{q_0}{k} \frac{1}{R_B - R_{HP}} \left( \frac{R_B^3}{9} - \frac{B_B^3}{4} + c_1 \ln(R_B) + c_2 \right) \\
  \label{eq:cyl_c2_expression2}
  \therefore c_2 &= T_B \frac{k}{q_0} (R_B - R_{HP}) + \frac{5}{36} R_B^3 - c_1 \ln(R_B)
\end{align}
Setting \eref{eq:cyl_c2_expression1} equal to \eref{eq:cyl_c2_expression2} yields an expression for $c_1$.
\begin{gather}
  T_{HP} \frac{k}{q_0} (R_B - R_{HP}) - \frac{R_{HP}^3}{9} + \frac{R_B \, R_{HP}^2}{4} - c_1 \ln(R_{HP}) = T_B \frac{k}{q_0} (R_B - R_{HP}) + \frac{5}{36} R_B^3 - c_1 \ln(R_B) \\
  (T_{HP} - T_{B}) \frac{k}{q_0} (R_B - R_{HP}) - \frac{R_{HP}^3}{9} + \frac{R_B \, R_{HP}^2}{4} - \frac{5}{36} R_B^3 = c_1 \left( \ln(R_{HP}) - \ln(R_B) \right) \\
  \label{eq:cyl_c1_expression}
  \therefore c_1 = \frac{1}{\ln\left(\frac{R_{HP}}{R_B}\right)} \left( (T_{HP} - T_B) \frac{k}{q_0} (R_B - R_{HP}) - \frac{R_{HP}^3}{9} + \frac{R_B \, R_{HP}^2}{4} - \frac{5}{36} R_B^3 \right)
\end{gather}

There is no ``clean'' final form of the solution for this problem.
The solution is given by the combination of \eref{eq:cyl_general}, \eref{eq:cyl_c2_expression2}, and \eref{eq:cyl_c1_expression}.
Evaluating the solution with the constants in \tref{tab:radial} yields the result plotted in \fref{fig:radial}.
Note that the temperature is treated as constant within the heat pipe.

\begin{table}
  \centering
  \caption{One-Dimension Cylindrical Heat Conduction Constants.}
  \label{tab:radial}
  \begin{tabular}{cc}
    \\
    \toprule
    Variable & Value \\
    \midrule
    $T_{HP}$ & 400 \\
    $T_{B}$  & 300 \\
    $R_{HP}$ & 0.5 \\
    $R_{B}$  & 2   \\
    $k$      & 0.1 \\
    $q_0$    & 100 \\
    \bottomrule
  \end{tabular}
\end{table}

\begin{figure}
  \centering
  \begin{subfigure}{0.4\textwidth}
    \centering
    \includegraphics[width=\columnwidth]{../figs/radial_line}
    \label{fig:radial_line}
  \end{subfigure}
  \begin{subfigure}{0.4\textwidth}
    \centering
    \includegraphics[width=\columnwidth]{../figs/radial_circle}
    \label{fig:radial_circle}
  \end{subfigure}
  \caption{One-Dimension Cylindrical Heat Conduction Solution.}
  \label{fig:radial}
\end{figure}

\section{One-Dimension Cartesian Heat Conduction with Variable Thermal Conductivity}

In realistic materials, thermal conductivity decreases as temperature increases.
This poses a challenge with solving the heat conduction equation as the thermal conductivity is dependent on the solution to the heat conduction equation.
Typically, thermal conductivities are weakly dependent on temperature so a simple fixed-point iteration is acceptable to treat the temperature dependence.
That is, the heat conduction equation is solved using an estimate for the material temperature and associated thermal conductivity and the material temperature is recalculated until convergence.

As the temperature dependence of thermal conductivity is an important effect in realistic materials, this analytic solution may serve as a benchmark for verifying computer programs designed to model heat conduction in realistic materials.
In one-dimensional Cartesian geometry, the heat conduction equation can be written
\begin{equation}
  \label{eq:var_cond_equation}
  - \frac{d}{dx} \left(k(x, T(x)) \frac{d}{dx} T(x)\right) = q(x).
\end{equation}
This solution will be restricted to the domain $x \in [0, L]$ and will assume the thermal conductivity to have the form
\begin{equation}
  \label{eq:var_cond_conductivity}
  k(x, T(x)) = k_0 - \alpha T(x)
\end{equation}
with a heat source of the form
\begin{equation}
  \label{eq:var_cond_source}
  q(x) = q_0 \cos \left( \frac{\pi \, x}{2L} \right).
\end{equation}
Boundary conditions will be treated as fixed temperatures such that
\begin{align}
  \label{eq:var_cond_T0}
  T(0) &= T_0 \\
  \label{eq:var_cond_TL}
  T(L) &= T_L.
\end{align}

The solution will begin by substituting the heat source definition from \eref{eq:var_cond_source} into the heat conduction equation from \eref{eq:var_cond_equation}.
\begin{equation}
  \label{eq:var_cond_1}
  - \frac{d}{dx} \left( k(x, T(x)) \frac{d}{dx} T(x) \right) = q_0 \cos \left( \frac{\pi \, x}{2L} \right)
\end{equation}
Then, \eref{eq:var_cond_1} is integrated with respect to $x$ as
\begin{equation}
  \label{eq:var_cond_2}
  -k(x, T(x)) \frac{d}{dx} T(x) = q_0 \frac{2L}{\pi} \sin \left( \frac{\pi \, x}{2L} \right) + c_1
\end{equation}
where $c_1$ is a constant of integration.
Then, the form of the variable thermal conductivity from \eref{eq:var_cond_conductivity} is substituted into \eref{eq:var_cond_2}.
\begin{align}
  (\alpha T(x) - k_0) \frac{d}{dx} T(x) &= q_0 \frac{2L}{\pi} \sin \left( \frac{\pi \, x}{2L} \right) + c_1 \\
  \label{eq:var_cond_3}
  \alpha T(x) \frac{d}{dx} T(x) - k_0 \frac{d}{dx} T(x) &= q_0 \frac{2L}{\pi} \sin \left( \frac{\pi \, x}{2L} \right) + c_1
\end{align}
Again, \eref{eq:var_cond_3} is integrated with respect to $x$.
\begin{align}
  \int \left( \alpha T(x) \frac{d}{dx} T(x) - k_0 \frac{d}{dx} T(x) \right) \; dx     &= \int \left( q_0 \frac{2L}{\pi} \sin \left( \frac{\pi \, x}{2L} \right) + c_1 \right) \; dx \\
  \label{eq:var_cond_4}
  \int \alpha T(x) \frac{d}{dx} T(x) \; dx - \int k_0 \frac{d}{dx} T(x) \; dx &= \int \left( q_0 \frac{2L}{\pi} \sin \left( \frac{\pi \, x}{2L} \right) + c_1 \right) \; dx
\end{align}
The first integral in the left-hand-side of \eref{eq:var_cond_4} can be evaluated using integration-by-parts and the second integral can be evaluated using the fundamental theory of calculus.
The resulting equation is
\begin{equation}
  \label{eq:var_cond_5}
  -\frac{\alpha}{2} T^2(x) + k_0 T(x) = q_0 \left( \frac{2L}{\pi} \right)^2 \cos \left( \frac{\pi \, x}{2L} \right) + c_1 x + c_2
\end{equation}
where $c_2$ is a constant of integration.
The quadratic equation can be used to evaluate \eref{eq:var_cond_5} to obtain a solution for $T(x)$.
\begin{equation}
  \label{eq:var_cond_sol}
  T(x) = \frac{1}{\alpha} \left( k_0 \pm \sqrt{k_0^2 - 2 \alpha \left( q_0 \left( \frac{2L}{\pi} \right)^2 \cos \left( \frac{\pi \, x}{2L} \right) + c_1 x + c_2 \right)} \right)
\end{equation}
For realistic calculations, the requirement that $T(x) > 0$ can be imposed in which case, the $\pm$ in \eref{eq:var_cond_sol} should be replaced with a minus-sign.

The constants of integration can be obtained by evaluating the boundary conditions.
First, consider the boundary condition at $x=0$ from \eref{eq:var_cond_T0}.
\begin{align}
  T(0) &= T_0 \\
  &= \frac{1}{\alpha} \left( k_0 \pm \sqrt{k_0^2 - 2 \alpha \left( q_0 \left( \frac{2L}{\pi} \right)^2 + c_2 \right)}\right) \\
  \alpha T_0 &= k_0 \pm \sqrt{k_0^2 - 2 \alpha \left( q_0 \left( \frac{2L}{\pi} \right)^2 + c_2 \right)} \\
  \alpha T_0 - k_0 &= \pm \sqrt{k_0^2 - 2 \alpha \left( q_0 \left( \frac{2L}{\pi} \right)^2 + c_2 \right)} \\
  (\alpha T_0 - k_0)^2 &= k_0^2 - 2 \alpha \left( q_0 \left( \frac{2L}{\pi} \right)^2 + c_2 \right) \\
  k_0^2 - (\alpha T_0 - k_0)^2 &= 2 \alpha \left( q_0 \left( \frac{2L}{\pi} \right)^2 + c_2 \right) \\
  \frac{1}{2\alpha} (k_0^2 - (\alpha T_0 - k_0)^2) &= q_0 \left( \frac{2L}{\pi} \right)^2 + c_2 \\
  \label{eq:var_cond_c2}
  \therefore c_2 &= \frac{1}{2\alpha} (k_0^2 - (\alpha T_0 - k_0)^2) - q_0 \left( \frac{2L}{\pi} \right)^2
\end{align}
A similar procedure can be followed to show that
\begin{equation}
  \label{eq:var_cond_c1}
  c_1 = \frac{1}{L} \left( \frac{1}{2\alpha} (k_0^2 - (\alpha T_L - k_0)^2) + c_2 \right).
\end{equation}
Finally, the solution to this heat conduction problem with variable thermal conductivity is given by \eref{eq:var_cond_sol} with the constants of integration given by \eref{eq:var_cond_c2} and \eref{eq:var_cond_c1}.

Similar to the derivation from \sref{sec:oned}, an analytic solution can be obtained for this problem assuming that thermal conductivity is constant at some effective thermal conductivity $\kbar$.
\begin{equation}
  \that(x) = \frac{q_0}{\kbar} \left( \frac{2L}{\pi} \right)^2 \left( \cos \left( \frac{\pi \, x}{2L} \right) + \frac{x}{L} - 1 \right) + \frac{T_L-T_0}{L} \, x + T_0
\end{equation}

Using the values in \tref{tab:var_cond}, an example solution is plotted in \fref{fig:var_cond}.
Numeric integration is used to obtain the effective thermal conductivity, $\kbar$, for this problem and the difference between the constant thermal conductivity and variable thermal conductivity solutions is also plotted in \fref{fig:linear}.

\begin{table}
  \centering
  \caption{One-Dimension Cartesian Heat Conduction with Variable Thermal Conductivity Constants.}
  \label{tab:var_cond}
  \begin{tabular}{cc}
    \toprule
    Variable & Value \\
    \midrule
    $T_0$ & $600$ \\
    $T_L$ & $300$ \\
    $L$   & $5$ \\
    $q_0$ & $200$ \\
    $k_0$ & $10$ \\
    $\alpha$ & $0.01$ \\
    $\kbar$ & $5.100360$ \\
    \bottomrule
    \end{tabular}
\end{table}

\begin{figure}
  \centering
  \begin{subfigure}{0.45\textwidth}
    \includegraphics[width=\textwidth]{../figs/var_cond_temperature}
    \caption{Temperature.}
    \label{fig:var_cond_temperature}
  \end{subfigure}
  \begin{subfigure}{0.45\textwidth}
    \includegraphics[width=\textwidth]{../figs/var_cond_diff}
    \caption{Temperature Difference.}
    \label{fig:var_cond_diff}
  \end{subfigure}
  \begin{subfigure}{0.45\textwidth}
    \includegraphics[width=\textwidth]{../figs/var_cond_conductivity}
    \caption{Thermal Conductivity.}
    \label{fig:var_cond_conductivity}
  \end{subfigure}
  \caption{One-Dimension Cartesian Heat Conduction with Variable Thermal Conductivity Constants.}
  \label{fig:var_cond}
\end{figure}

\end{document}
